\chapter{Future Work}
\section{Methodological Enhancements}
The extension compared three DE variants on one representation; future experiments could test more adaptive controls, other search methods, or more engineered inputs using a protocol fixed before evaluation. Population size, mutation weight and crossover rate should be varied systematically with equal budgets. Such experiments must be presented as new work rather than retroactively improving the current frozen result.
\section{Data and Validation}
Repeated outer group splits or nested group cross-validation would assess sensitivity to partition choice. External concrete datasets would test transfer beyond the current measurement collection. Age-stratified errors, uncertainty intervals and ingredient-level interpretation could add useful context if designed and validated carefully.

Because the current holdout has already been evaluated, it should not become a feedback channel for repeated improvement. New model-selection decisions require a clearly separated evaluation design. Feature importance and ablation studies should report their methods and uncertainty rather than infer rankings from a single selected configuration.
\section{Sustainable Development Goals}
Reliable materials modeling could be relevant to SDG 9 (Industry, Innovation and Infrastructure) and SDG 12 (Responsible Consumption and Production). This is a potential connection, not a measured sustainability outcome. The present experiment does not quantify material savings, carbon reduction or safety improvements. A future constrained mixture-design study would need environmental data and engineering validation to assess those outcomes.
